%
En kasse med masse $m = \SI{20}{\kilo\gram}$ sklir bortover et gulv med startfart $v_0 = \SI{3.0}{\metre\per\second}$. Den eneste kraften som bremser kassen, er friksjonskraften $R = \SI{40}{\newton}$.
\begin{enumerate}[a)]
    \item Hvor stor er endringen i kinetisk energi fra kassen begynner å skli til den står stille?
    \item Bruk arbeid-energi-setningen til å finne hvor langt kassen sklir før den stopper.
    \item Hvor langt ville kassen skli dersom startfarten var dobbelt så stor?
\end{enumerate}
